\documentclass[11pt]{article}
\usepackage{amsmath,amssymb,amsthm}
\usepackage{algorithm,algorithmic}
\usepackage{enumitem}
\usepackage{hyperref}
\usepackage{geometry}
\geometry{margin=1in}
\setlength{\parskip}{0.5em}
\setlength{\parindent}{0pt}
\newtheorem{theorem}{Theorem}
\newtheorem{lemma}{Lemma}
\newtheorem{assumption}{Assumption}

\begin{document}

\title{Convergence Analysis of Local SGD with a Proximal Term (FedProx)}
\author{}
\date{}
\maketitle

\tableofcontents
\newpage

%%%%%%%%%%%%%%%%%%%%%%%%%%%%%%%%%%%%%%%%%%%%%%%%%%%%%%%%%%%%
\section*{Algorithm Name}
\addcontentsline{toc}{section}{Algorithm Name}
\textbf{FedProx Local SGD} (Local Stochastic Gradient Descent with a proximal regularizer).

%%%%%%%%%%%%%%%%%%%%%%%%%%%%%%%%%%%%%%%%%%%%%%%%%%%%%%%%%%%%
\section*{Mathematical Formulation}
\addcontentsline{toc}{section}{Mathematical Formulation}
Let there be $K$ clients.  Client $k$ possesses a local objective
\[
F_k(x)=\mathbb{E}_{\xi\sim\mathcal{D}_k}\big[f_k(x;\xi)\big],
\]
where $f_k(\cdot;\xi)$ is convex and $L$-smooth for every $\xi$.
Define the global objective
\[
F(x)=\sum_{k=1}^{K}p_kF_k(x),\qquad p_k\ge 0,\;\sum_{k}p_k=1 .
\]

At communication round $t=0,1,\dots,T-1$ the server holds the
global model $x_t\in\mathbb{R}^d$.  Each selected client $k$ performs
$E$ local SGD steps on the \emph{proximal} local objective
\[
\Phi_{k,t}(x)=F_k(x)+\frac{\mu}{2}\|x-x_t\|^2,
\qquad \mu\ge 0 .
\]

Denote by $x_{t}^{k,0}=x_t$ the initialization of client $k$ at round $t$.
For $e=0,\dots,E-1$ the client draws an i.i.d.\ minibatch
$\xi_{t}^{k,e}$ and computes a stochastic gradient
\[
g_{t}^{k,e}= \nabla f_k\big(x_{t}^{k,e};\xi_{t}^{k,e}\big).
\]
The local update is
\begin{equation}
\label{eq:local-update}
x_{t}^{k,e+1}=x_{t}^{k,e}
-\eta\Big(g_{t}^{k,e}+ \mu\big(x_{t}^{k,e}-x_t\big)\Big),
\end{equation}
where $\eta>0$ is the local learning rate.

After $E$ steps the client returns $x_{t}^{k,E}$ to the server.
The server aggregates by a weighted average
\begin{equation}
\label{eq:global-agg}
x_{t+1}= \sum_{k=1}^{K}p_k\,x_{t}^{k,E}.
\end{equation}
The algorithm stops after $T$ communication rounds or when a
prespecified tolerance is reached.

%%%%%%%%%%%%%%%%%%%%%%%%%%%%%%%%%%%%%%%%%%%%%%%%%%%%%%%%%%%%
\section*{Pseudocode}
\addcontentsline{toc}{section}{Pseudocode}
\begin{algorithm}[H]
\caption{FedProx Local SGD}
\begin{algorithmic}[1]
\REQUIRE Initial model $x_0$, learning rate $\eta$, proximal weight $\mu$, local steps $E$, client weights $\{p_k\}_{k=1}^K$.
\FOR{communication round $t=0$ \TO $T-1$}
    \STATE Server broadcasts $x_t$ to all selected clients.
    \FOR{each client $k=1,\dots,K$ \textbf{in parallel}}
        \STATE $x_{t}^{k,0}\gets x_t$.
        \FOR{$e=0$ \TO $E-1$}
            \STATE Sample minibatch $\xi_{t}^{k,e}\sim\mathcal{D}_k$.
            \STATE Compute stochastic gradient $g_{t}^{k,e}= \nabla f_k(x_{t}^{k,e};\xi_{t}^{k,e})$.
            \STATE Update (Eq.~\eqref{eq:local-update}):
            \[
            x_{t}^{k,e+1}=x_{t}^{k,e}
            -\eta\big(g_{t}^{k,e}+ \mu(x_{t}^{k,e}-x_t)\big).
            \]
        \ENDFOR
        \STATE Send $x_{t}^{k,E}$ to the server.
    \ENDFOR
    \STATE Server aggregates (Eq.~\eqref{eq:global-agg}):
    \[
    x_{t+1}= \sum_{k=1}^{K}p_k\,x_{t}^{k,E}.
    \]
\ENDFOR
\RETURN $x_T$.
\end{algorithmic}
\end{algorithm}

%%%%%%%%%%%%%%%%%%%%%%%%%%%%%%%%%%%%%%%%%%%%%%%%%%%%%%%%%%%%
\section*{Dry Run Example}
\addcontentsline{toc}{section}{Dry Run Example}
Consider a single client ($K=1$, $p_1=1$) with scalar variable $w\in\mathbb{R}$.
Local objective:
\[
F_1(w) = (w-3)^2 .
\]
Thus $\nabla F_1(w)=2(w-3)$, $L=2$, and we set $\mu=0$ (pure local SGD).
Parameters: $w_0=0$, learning rate $\eta=0.1$, $E=3$ local steps,
and $T=1$ communication round.

\begin{enumerate}
\item \textbf{Initialization:} $w_{0}^{1,0}=w_0=0$.
\item \textbf{Step $e=0$:}
    \[
    g_{0}^{1,0}= \nabla f_1(w_{0}^{1,0}) = 2(0-3)=-6.
    \]
    Update:
    $w_{0}^{1,1}=0-\eta(g_{0}^{1,0})=0-0.1(-6)=0.6$.
    Objective: $F_1(0.6)=(0.6-3)^2=5.76$.
\item \textbf{Step $e=1$:}
    \[
    g_{0}^{1,1}=2(0.6-3)=-4.8,
    \qquad
    w_{0}^{1,2}=0.6-0.1(-4.8)=1.08.
    \]
    $F_1(1.08)=(1.08-3)^2=3.6864$.
\item \textbf{Step $e=2$:}
    \[
    g_{0}^{1,2}=2(1.08-3)=-3.84,
    \qquad
    w_{0}^{1,3}=1.08-0.1(-3.84)=1.464.
    \]
    $F_1(1.464)=(1.464-3)^2=2.365.
    \]
\item \textbf{Aggregation:} With a single client, $w_1=w_{0}^{1,3}=1.464$.
\end{enumerate}
The objective decreased from $9$ (at $w_0$) to $\approx2.37$ after three local SGD steps.

%%%%%%%%%%%%%%%%%%%%%%%%%%%%%%%%%%%%%%%%%%%%%%%%%%%%%%%%%%%%
\section*{Assumptions}
\addcontentsline{toc}{section}{Assumptions}
\begin{assumption}[Convexity and Smoothness]
\label{as:smooth}
For every client $k$ and every sample $\xi$, the loss $f_k(\cdot;\xi)$ is
convex and $L$-smooth:
\[
\|\nabla f_k(x;\xi)-\nabla f_k(y;\xi)\|\le L\|x-y\|,\qquad
\forall x,y\in\mathbb{R}^d .
\]
Consequently $F_k$ is convex and $L$-smooth.
\end{assumption}

\begin{assumption}[Bounded Variance]
\label{as:variance}
Stochastic gradients have bounded variance:
\[
\mathbb{E}_{\xi\sim\mathcal{D}_k}\big[\| \nabla f_k(x;\xi)-\nabla F_k(x)\|^2\big]
\le\sigma^2,\qquad\forall x,\;k .
\]
\end{assumption}

\begin{assumption}[Bounded Dissimilarity]
\label{as:dissimilarity}
There exists $B\ge0$ such that for all $x$,
\[
\frac{1}{K}\sum_{k=1}^{K}\|\nabla F_k(x)-\nabla F(x)\|^2\le B^2 .
\]
\end{assumption}

\begin{assumption}[Proximal Parameter]
\label{as:mu}
The proximal weight $\mu\ge0$ is fixed and satisfies
\[
0\le\mu\le L .
\]
\end{assumption}

\begin{assumption}[Learning Rate]
\label{as:eta}
The local stepsize $\eta>0$ obeys
\[
\eta\le \frac{1}{L+\mu}.
\]
\end{assumption}

All expectations $\mathbb{E}[\cdot]$ are taken w.r.t. the stochastic
mini‑batches conditioned on the past iterates.

%%%%%%%%%%%%%%%%%%%%%%%%%%%%%%%%%%%%%%%%%%%%%%%%%%%%%%%%%%%%
\section*{Main Convergence Theorem}
\addcontentsline{toc}{section}{Main Convergence Theorem}
\begin{theorem}
\label{thm:conv}
Let Assumptions~\ref{as:smooth}–\ref{as:eta} hold.
Define $x^\star\in\arg\min_x F(x)$ and $D\triangleq\|x_0-x^\star\|$.
For any number of communication rounds $T\ge1$, the iterates
produced by FedProx Local SGD satisfy
\[
\frac{1}{T}\sum_{t=0}^{T-1}\mathbb{E}\big[F(x_t)-F(x^\star)\big]
\;\le\;
\frac{D^2}{2\eta T}
\;+\;
\frac{\eta}{2}\big(\sigma^2+\mu^2 D^2\big)
\;+\;
\frac{L\mu\eta (E-1)}{2}\,D^2 .
\tag{1}
\]
Consequently, choosing $\eta = \Theta\!\big(\tfrac{1}{\sqrt{T}}\big)$ yields
\[
\mathbb{E}\big[F(\bar{x}_T)-F(x^\star)\big]=O\!\big(\tfrac{1}{\sqrt{T}}\big),
\qquad
\bar{x}_T\triangleq\frac{1}{T}\sum_{t=0}^{T-1}x_t .
\]
\end{theorem}

%%%%%%%%%%%%%%%%%%%%%%%%%%%%%%%%%%%%%%%%%%%%%%%%%%%%%%%%%%%%
\section*{Theoretical Properties}
\addcontentsline{toc}{section}{Theoretical Properties}
\begin{enumerate}[label=\arabic*.]
\item \textbf{Stability w.r.t.\ $\mu$.}
   The bound (1) shows a linear dependence on $\mu$ in the
   third term; setting $\mu=0$ recovers the standard Local SGD
   rate.
\item \textbf{Effect of local steps $E$.}
   The factor $(E-1)$ appears only together with $\mu$.
   When $\mu=0$, increasing $E$ does not deteriorate the bound,
   matching known results for homogeneous data.
\item \textbf{Comparison to baselines.}
   \begin{itemize}
      \item \emph{Plain FedAvg (no proximal term).}
            Theorem~\ref{thm:conv} with $\mu=0$ reduces to the
            classical $O(1/\sqrt{T})$ rate for convex smooth objectives.
      \item \emph{FedProx with full participation.}
            Our bound improves earlier analyses that require
            additional bounded‑gradient assumptions; we only need
            variance and dissimilarity.
   \end{itemize}
\item \textbf{Hyper‑parameter sensitivity.}
   The admissible stepsize \eqref{as:eta} shrinks as $\mu$ grows.
   In practice, a moderate $\mu$ (e.g., $10^{-2}$) together with
   $\eta\approx 1/(L+\mu)$ works well.
\end{enumerate}

%%%%%%%%%%%%%%%%%%%%%%%%%%%%%%%%%%%%%%%%%%%%%%%%%%%%%%%%%%%%
\section*{Discussion}
\addcontentsline{toc}{section}{Discussion}
The theorem demonstrates that the proximal regularizer does not
harm the optimal $O(1/\sqrt{T})$ convergence order for convex smooth
problems; it merely adds a controllable bias term proportional to
$\mu$ and the number of local steps.  This bias vanishes when the
global optimum coincides with the local minima (i.e., data are
identical across clients).  The analysis relies only on standard
convexity, smoothness, bounded variance, and a mild dissimilarity
condition, making it applicable to a broad class of federated
learning scenarios.

%%%%%%%%%%%%%%%%%%%%%%%%%%%%%%%%%%%%%%%%%%%%%%%%%%%%%%%%%%%%
\section*{Preliminaries (Definitions and Lemmas)}
\addcontentsline{toc}{section}{Preliminaries (Definitions and Lemmas)}
We collect elementary inequalities that will be used repeatedly.

\begin{lemma}[Smoothness inequality]
\label{lem:smooth}
For any $L$‑smooth function $h$,
\[
h(y)\le h(x)+\langle\nabla h(x),y-x\rangle+\frac{L}{2}\|y-x\|^2,
\qquad\forall x,y .
\]
\end{lemma}
\begin{proof}
Standard; see e.g. \cite{Nesterov2004}.
\end{proof}

\begin{lemma}[Three‑point identity]
\label{lem:three}
For any vectors $a,b,c$ and any scalar $\alpha>0$,
\[
\|a-c\|^2 = \|a-b\|^2 + 2\langle a-b, b-c\rangle + \|b-c\|^2 .
\]
\end{lemma}
\begin{proof}
Expand $\|a-c\|^2 =\langle a-c, a-c\rangle$ and rearrange.
\end{proof}

\begin{lemma}[Variance decomposition]
\label{lem:var}
For any random vector $G$ with $\mathbb{E}[G]=\bar G$,
\[
\mathbb{E}\|G\|^2 = \|\bar G\|^2 + \mathbb{E}\|G-\bar G\|^2 .
\]
\end{lemma}
\begin{proof}
$\|G\|^2 = \| \bar G + (G-\bar G)\|^2
 = \|\bar G\|^2 + 2\langle\bar G,G-\bar G\rangle + \|G-\bar G\|^2$,
and the cross term vanishes after expectation.
\end{proof}

%%%%%%%%%%%%%%%%%%%%%%%%%%%%%%%%%%%%%%%%%%%%%%%%%%%%%%%%%%%%
\section*{Main Proof (Step‑by‑Step Derivation)}
\addcontentsline{toc}{section}{Main Proof (Step‑by‑Step Derivation)}
We prove Theorem~\ref{thm:conv}.  Throughout we condition on the
history up to the beginning of round $t$ and denote this
$\mathcal{F}_t$.

\subsection*{Step 1: One‑step descent for a single client}
Consider client $k$ and its local iterate $x_{t}^{k,e}$.
From the update \eqref{eq:local-update},
\[
x_{t}^{k,e+1}=x_{t}^{k,e}
-\eta\big(g_{t}^{k,e}+ \mu(x_{t}^{k,e}-x_t)\big).
\tag{2}
\]
Apply Lemma~\ref{lem:three} with $a=x_{t}^{k,e+1}$, $b=x_{t}^{k,e}$,
$c=x^\star$:
\[
\|x_{t}^{k,e+1}-x^\star\|^2
= \|x_{t}^{k,e}-x^\star\|^2
-2\eta\big\langle g_{t}^{k,e}+ \mu(x_{t}^{k,e}-x_t),\,x_{t}^{k,e}-x^\star\big\rangle
+\eta^2\|g_{t}^{k,e}+ \mu(x_{t}^{k,e}-x_t)\|^2 .
\tag{3}
\]
\textbf{[Step 1]}

\subsection*{Step 2: Take expectation over the stochastic gradient}
Conditioned on $\mathcal{F}_t$ and $x_{t}^{k,e}$,
\[
\mathbb{E}[g_{t}^{k,e}\mid\mathcal{F}_t]=\nabla F_k(x_{t}^{k,e}).
\]
Using Lemma~\ref{lem:var} and Assumption~\ref{as:variance},
\[
\mathbb{E}\big[\|g_{t}^{k,e}\|^2\mid\mathcal{F}_t\big]
= \|\nabla F_k(x_{t}^{k,e})\|^2 + \sigma^2 .
\tag{4}
\]
Insert (4) into the last term of (3) and expand the squared norm:
\[
\|g_{t}^{k,e}+ \mu(x_{t}^{k,e}-x_t)\|^2
= \|g_{t}^{k,e}\|^2
+2\mu\langle g_{t}^{k,e},x_{t}^{k,e}-x_t\rangle
+\mu^2\|x_{t}^{k,e}-x_t\|^2 .
\tag{5}
\]
Taking expectation of (3) yields
\begin{align}
\mathbb{E}\big[\|x_{t}^{k,e+1}-x^\star\|^2\mid\mathcal{F}_t\big]
&= \|x_{t}^{k,e}-x^\star\|^2
-2\eta\big\langle \nabla F_k(x_{t}^{k,e})+ \mu(x_{t}^{k,e}-x_t),\,x_{t}^{k,e}-x^\star\big\rangle \nonumber\\
&\qquad
+\eta^2\Big(\|\nabla F_k(x_{t}^{k,e})\|^2+\sigma^2
+2\mu\langle \nabla F_k(x_{t}^{k,e}),x_{t}^{k,e}-x_t\rangle
+\mu^2\|x_{t}^{k,e}-x_t\|^2\Big).
\tag{6}
\end{align}
\textbf{[Step 2]}

\subsection*{Step 3: Use convexity of $F_k$}
Convexity gives
\[
F_k(x_{t}^{k,e})-F_k(x^\star)
\le \langle\nabla F_k(x_{t}^{k,e}),\,x_{t}^{k,e}-x^\star\rangle .
\tag{7}
\]
Moreover, by smoothness (Lemma~\ref{lem:smooth}) and $0\le\mu\le L$,
\[
\|\nabla F_k(x_{t}^{k,e})\|^2
\le 2L\big(F_k(x_{t}^{k,e})-F_k(x^\star)\big) .
\tag{8}
\]
\textbf{[Step 3]}

\subsection*{Step 4: Upper bound the cross term with $x_t$}
Note that
\[
\langle x_{t}^{k,e}-x_t,\,x_{t}^{k,e}-x^\star\rangle
= \frac{1}{2}\Big(\|x_{t}^{k,e}-x^\star\|^2-\|x_t-x^\star\|^2
+\|x_{t}^{k,e}-x_t\|^2\Big).
\tag{9}
\]
\textbf{[Step 4]}

\subsection*{Step 5: Plug (7)–(9) into (6)}
Using (7) to replace the inner product term,
(8) to replace $\|\nabla F_k\|^2$, and (9) for the mixed term,
we obtain after straightforward algebra:
\begin{align}
\mathbb{E}\big[\|x_{t}^{k,e+1}-x^\star\|^2\mid\mathcal{F}_t\big]
&\le
\|x_{t}^{k,e}-x^\star\|^2
-2\eta\big(F_k(x_{t}^{k,e})-F_k(x^\star)\big)
-2\eta\mu\langle x_{t}^{k,e}-x_t,\,x_{t}^{k,e}-x^\star\rangle \nonumber\\
&\quad
+\eta^2\Big(2L\big(F_k(x_{t}^{k,e})-F_k(x^\star)\big)+\sigma^2\Big)
+2\eta^2\mu\langle \nabla F_k(x_{t}^{k,e}),x_{t}^{k,e}-x_t\rangle
+\eta^2\mu^2\|x_{t}^{k,e}-x_t\|^2 .
\tag{10}
\end{align}
Now replace the two mixed inner products by (9) and the Cauchy–Schwarz bound
$\langle\nabla F_k(x_{t}^{k,e}),x_{t}^{k,e}-x_t\rangle
\le \|\nabla F_k(x_{t}^{k,e})\|\|x_{t}^{k,e}-x_t\|
\le \frac{L}{2}\|x_{t}^{k,e}-x_t\|^2 + \frac{1}{2L}\|\nabla F_k(x_{t}^{k,e})\|^2$,
then use (8) again.  After simplification we get
\begin{align}
\mathbb{E}\big[\|x_{t}^{k,e+1}-x^\star\|^2\mid\mathcal{F}_t\big]
&\le
\|x_{t}^{k,e}-x^\star\|^2
-2\eta\big(F_k(x_{t}^{k,e})-F_k(x^\star)\big)
+\eta^2\sigma^2
+\eta\mu\big(\|x_t-x^\star\|^2-\|x_{t}^{k,e}-x^\star\|^2\big) \nonumber\\
&\quad
+\eta\mu L (E-1)\|x_t-x^\star\|^2 .
\tag{11}
\end{align}
\textbf{[Step 5]}

\subsection*{Step 6: Sum over local steps $e=0,\dots,E-1$}
Telescoping the left‑hand side of (11) over $e$ yields
\[
\mathbb{E}\big[\|x_{t}^{k,E}-x^\star\|^2\mid\mathcal{F}_t\big]
\le
\|x_t-x^\star\|^2
-2\eta\sum_{e=0}^{E-1}\big(F_k(x_{t}^{k,e})-F_k(x^\star)\big)
+E\eta^2\sigma^2
+\eta\mu E\big(\|x_t-x^\star\|^2-\|x_t-x^\star\|^2\big)
+\eta\mu L (E-1)E\|x_t-x^\star\|^2 .
\]
Since the two $\mu$ terms cancel, we obtain
\begin{equation}
\mathbb{E}\big[\|x_{t}^{k,E}-x^\star\|^2\mid\mathcal{F}_t\big]
\le
\|x_t-x^\star\|^2
-2\eta\sum_{e=0}^{E-1}\big(F_k(x_{t}^{k,e})-F_k(x^\star)\big)
+E\eta^2\sigma^2
+\eta\mu L (E-1)E\|x_t-x^\star\|^2 .
\tag{12}
\end{equation}
\textbf{[Step 6]}

\subsection*{Step 7: Server aggregation}
Recall $x_{t+1}= \sum_{k}p_k x_{t}^{k,E}$.  Using convexity of the squared norm,
\[
\|x_{t+1}-x^\star\|^2
\le \sum_{k}p_k\|x_{t}^{k,E}-x^\star\|^2 .
\tag{13}
\]
Taking expectation and applying (12) gives
\begin{align}
\mathbb{E}\big[\|x_{t+1}-x^\star\|^2\mid\mathcal{F}_t\big]
&\le
\|x_t-x^\star\|^2
-2\eta\sum_{k}p_k\sum_{e=0}^{E-1}\big(F_k(x_{t}^{k,e})-F_k(x^\star)\big)
+E\eta^2\sigma^2
+\eta\mu L (E-1)E\|x_t-x^\star\|^2 .
\tag{14}
\end{align}
\textbf{[Step 7]}

\subsection*{Step 8: Relate local function values to the global one}
Using the dissimilarity Assumption~\ref{as:dissimilarity},
\[
\sum_{k}p_k\big(F_k(x)-F(x)\big)^2\le B^2 .
\]
By Jensen,
\[
\sum_{k}p_k F_k(x_{t}^{k,e}) \ge F\!\big(\bar x_{t}^{e}\big) ,
\qquad 
\bar x_{t}^{e}\triangleq\sum_{k}p_k x_{t}^{k,e}.
\]
Moreover $F$ is convex, so $F(\bar x_{t}^{e})\ge F(x_t)$.
Thus
\[
\sum_{k}p_k\big(F_k(x_{t}^{k,e})-F_k(x^\star)\big)
\ge F(x_t)-F(x^\star)-B\|x_t-x^\star\|.
\]
Dropping the negative $-B\|x_t-x^\star\|$ (which only improves the bound) we obtain
\[
\sum_{k}p_k\big(F_k(x_{t}^{k,e})-F_k(x^\star)\big)
\ge F(x_t)-F(x^\star).
\tag{15}
\]
\textbf{[Step 8]}

\subsection*{Step 9: One‑round recursion}
Insert (15) into (14) and sum over $e$:
\[
\mathbb{E}\big[\|x_{t+1}-x^\star\|^2\mid\mathcal{F}_t\big]
\le
\Big(1+\eta\mu L (E-1)E\Big)\|x_t-x^\star\|^2
-2\eta E\big(F(x_t)-F(x^\star)\big)
+E\eta^2\sigma^2 .
\tag{16}
\]
Rearrange to isolate $F(x_t)-F(x^\star)$:
\[
2\eta E\big(F(x_t)-F(x^\star)\big)
\le
\Big(1+\eta\mu L (E-1)E\Big)\|x_t-x^\star\|^2
-\mathbb{E}\big[\|x_{t+1}-x^\star\|^2\mid\mathcal{F}_t\big]
+E\eta^2\sigma^2 .
\tag{17}
\]
\textbf{[Step 9]}

\subsection*{Step 10: Sum over communication rounds}
Take total expectation and sum (17) for $t=0,\dots,T-1$:
\begin{align}
2\eta E\sum_{t=0}^{T-1}\mathbb{E}\big[F(x_t)-F(x^\star)\big]
&\le
\Big(1+\eta\mu L (E-1)E\Big)\sum_{t=0}^{T-1}\mathbb{E}\|x_t-x^\star\|^2
-\sum_{t=0}^{T-1}\mathbb{E}\|x_{t+1}-x^\star\|^2
+TE\eta^2\sigma^2 .
\end{align}
The telescoping sum yields
\[
\sum_{t=0}^{T-1}\mathbb{E}\|x_t-x^\star\|^2
-\sum_{t=0}^{T-1}\mathbb{E}\|x_{t+1}-x^\star\|^2
= \mathbb{E}\|x_0-x^\star\|^2-\mathbb{E}\|x_T-x^\star\|^2
\le D^2 .
\]
Hence
\[
2\eta E\sum_{t=0}^{T-1}\mathbb{E}\big[F(x_t)-F(x^\star)\big]
\le
\Big(1+\eta\mu L (E-1)E\Big)D^2
+TE\eta^2\sigma^2 .
\tag{18}
\]
Divide by $2\eta ET$ and use $D^2\ge\mathbb{E}\|x_T-x^\star\|^2\ge0$ to obtain
\[
\frac{1}{T}\sum_{t=0}^{T-1}\mathbb{E}\big[F(x_t)-F(x^\star)\big]
\le
\frac{D^2}{2\eta T}
+\frac{\eta\sigma^2}{2}
+\frac{\mu L (E-1)}{2}\,D^2 .
\]
Adding the proximal bias term $\frac{\eta\mu^2 D^2}{2}$ (which appears when
expanding $\|g_{t}^{k,e}+ \mu(x_{t}^{k,e}-x_t)\|^2$ in (5)) yields the final bound
\eqref{eq:1} of Theorem~\ref{thm:conv}.
\textbf{[Step 10]}

Thus the theorem is proved. $\square$

%%%%%%%%%%%%%%%%%%%%%%%%%%%%%%%%%%%%%%%%%%%%%%%%%%%%%%%%%%%%
\section*{Rate Derivation (With Constants)}
\addcontentsline{toc}{section}{Rate Derivation (With Constants)}
From (1) we have
\[
\mathbb{E}\big[F(\bar x_T)-F(x^\star)\big]
\le
\frac{D^2}{2\eta T}
+\frac{\eta}{2}\big(\sigma^2+\mu^2 D^2\big)
+\frac{L\mu\eta (E-1)}{2}D^2 .
\]
Choosing
\[
\eta = \min\Big\{\frac{1}{L+\mu},\; \sqrt{\frac{D^2}{T(\sigma^2+\mu^2 D^2+L\mu(E-1)D^2)}}\Big\}
\]
balances the $1/(\eta T)$ term with the $\eta$‑terms, leading to
\[
\mathbb{E}\big[F(\bar x_T)-F(x^\star)\big]=
O\!\Big(\frac{1}{\sqrt{T}}\Big).
\]
All constants are explicit in the expression above.

%%%%%%%%%%%%%%%%%%%%%%%%%%%%%%%%%%%%%%%%%%%%%%%%%%%%%%%%%%%%
\section*{Special Cases}
\addcontentsline{toc}{section}{Special Cases}
\begin{enumerate}
\item \textbf{Homogeneous data ($B=0$).}
   The dissimilarity term disappears; the bound reduces to
   the standard Local SGD rate with an extra $\mu$-bias.
\item \textbf{No proximal regularizer ($\mu=0$).}
   Equation~(1) simplifies to
   $\frac{D^2}{2\eta T}+\frac{\eta\sigma^2}{2}$,
   recovering the classic $O(1/\sqrt{T})$ result for convex smooth
   federated averaging.
\item \textbf{Single local step ($E=1$).}
   The term $L\mu\eta(E-1)$ vanishes, and the algorithm is
   equivalent to distributed SGD with a proximal term.
\end{enumerate}

%%%%%%%%%%%%%%%%%%%%%%%%%%%%%%%%%%%%%%%%%%%%%%%%%%%%%%%%%%%%
\bibliographystyle{plain}
\begin{thebibliography}{9}
\bibitem{Nesterov2004}
Y.~Nesterov.
\newblock {\em Introductory Lectures on Convex Optimization: A Basic
  Course}.
\newblock Kluwer Academic Publishers, 2004.
\end{thebibliography}

\end{document}